\documentclass[10pt,a4paper]{amsart}
\usepackage[margin=19mm]{geometry}
\usepackage{amsmath,amssymb,mathtools}
\usepackage[scr=rsfs,cal=euler]{mathalfa}
\usepackage{xfrac}
\usepackage[unicode,hidelinks,bookmarks=false]{hyperref}
\newcommand{\ie}{i.e.\ }
\newcommand{\cf}{cf.\ }
\newcommand{\resp}{respectively\ }
\newcommand{\Subsection}{\subsection}
\providecommand{\nobreakdash}{\nobreak\hbox{-}}
\setlength{\emergencystretch}{2em}
\multlinegap0pt
\usepackage{mathtools}
\usepackage{amssymb}
\usepackage{bm}
\usepackage{dsfont}
\usepackage{xcolor}
\newtheorem{theorem}{Theorem}
\newtheorem{corollary}{Corollary}
\def\N{{\mathbb N}}
\def\R{{\mathbb R}}
\def\Z{{\mathbb Z}}
\newcommand{\ft}{{\mathcal{F}}}
\newcommand{\ift}{{\mathcal{F}^{-1}}}
\newcommand{\supp}{{\mathrm{supp}}}
\title{Nonlinear determination and phase retrieval under unimodular constraints}
\author{Lukas Liehr, Tomasz Szczepanski}
\date{Source: arXiv:2601.00403v1 (CC BY 4.0)}
\begin{document}
\maketitle

\noindent\emph{Source and licence.} Nonlinear determination and phase retrieval under unimodular constraints. Version arXiv:2601.00403v1 (CC BY 4.0), distributed by arXiv under the Creative Commons Attribution 4.0 International licence, http://creativecommons.org/licenses/by/4.0/, as recorded in the arXiv licence metadata for this exact version and shown on the abstract page https://arxiv.org/abs/2601.00403v1. Full licence notice: \texttt{licenses/arxiv-2601.00403-CC-BY-4.0.txt}.\par\smallskip

\noindent\textit{Editorial scope.} Counted unit: the article's theorem thm:lattice\_density\_characterization (source0.tex lines 181--183). The article's complete original proof of it is delivered with it (source0.tex lines 689--774, one \texttt{proof} environment of 86 lines, which the article places away from the statement). No mathematics is written, repaired or re-derived: the statement and the proof are the article's own lines, in the article's own notation, and no retained line is substituted or re-flowed.  Retained uncounted as the source-local context the proof needs: lines 662--669.  Nothing was omitted from any retained span, so the package carries no omission record.  No retained line contains a citation, so the wrapper carries no bibliography and no bibliography entry.  No retained line contains a figure environment, an image-inclusion command or drawing code, so no image is delivered and the package declares no asset dependency.  Editorial formatting only, in addition: an AMS \texttt{amsart} preamble that restates the article's own declarations and macros used by the retained lines, namely the environments \texttt{theorem}, \texttt{corollary} on separate counters, together with the standard packages those lines need.  The retained environments print in this document as Theorem 1, Corollary 1 in order of appearance, whereas the article prints the same items with the numbering of its own full document.  The complete native source of the article as distributed in the arXiv e-print for this exact version is bundled unchanged as \texttt{sources/191997/source0.tex}.
\begin{corollary}\label{cor:PW_recurrence}
Let $\Theta = \{ e^{2\pi i k/n} : k=0,\dots,n-1 \}$, let $\Lambda \subseteq \R$ and let $a>0$. Then $E(\Lambda)$ fails $\Theta$-PR in $L^2[0,a]$ if and only if there exist a cover $\{\Lambda_j\}_{j \in \N}$ of $\Lambda$ and functions $\{ x_j \}_{j \in \N} \subseteq PW_a \setminus \{ 0 \}$ satisfying
\begin{enumerate}
        \item $x_j$ vanishes on $\Lambda_j$,
        \item $x_1$ and $x_2$ are linearly independent,
        \item for every $j \geq 3$, it holds that $x_j = (1+e^{\frac{2\pi i }{n}})x_{j-1} - e^{\frac{2\pi i }{n}} x_{j-2}$.
    \end{enumerate}
\end{corollary}

\begin{theorem}\label{thm:lattice_density_characterization}
    Let $\Lambda \subseteq \R$ be a lattice and let $\Theta = \{ e^{2\pi i k/n} : k=0,\dots,n-1 \}$. Then $E(\Lambda)$ does $\Theta$-PR in $L^2[0,a]$ if and only if $D(\Lambda) \geq na$.
\end{theorem}

\begin{proof}[Proof of Theorem \ref{thm:lattice_density_characterization}]
Without loss of generality, we can assume that $a=1$. The general result follows by scaling. Moreover, we can assume that $n \geq 2$, as the case $n=1$ is a consequence of the classical Fourier uniqueness theorem.

\textbf{Sufficiency.}
For $\lambda \in \Lambda$ define as usual $e_\lambda(x) = e^{2\pi i \lambda x}$. Let $f,h \in L^2[0,1]$ and suppose that for every $\lambda \in \Lambda$ there exists $\theta_\lambda \in \Theta$ such that
\begin{equation}\label{eq:eeeee}
    \langle f,e_\lambda \rangle = \theta_\lambda \langle h,e_\lambda \rangle.
\end{equation}
Using the definition of the Fourier transform it follows that for $\lambda \in \Lambda$, the identity \eqref{eq:eeeee} is equivalent to the property that
$$
\ft(f-\theta_\lambda h)(\lambda) =0.
$$
Now let $\omega_j=e^{2\pi i j/n}$ and define for $j \in \{ 0, \dots, n-1 \}$ functions $p_j \in L^2[0,1]$ via
\begin{equation}\label{eq:prod_analytic}
        p_j \coloneqq  f - \omega_j h.
    \end{equation}
Consider each $p_j$ as a function in $L^2(\R)$ with support in $[0,1]$. By our assumption on $f,h$ we have that
$$
\prod_{j=0}^{n-1} \ft(p_j)(\lambda) = 0, \quad \lambda \in \Lambda.
$$
Using the convolution theorem for the Fourier transform, the latter relation is equivalent to the property that the $n$-fold convolution
$$
p \coloneqq p_1 * \dots * p_n
$$
satisfies
$$
\ft p(\lambda) = \langle p, e_\lambda \rangle =0, \quad \lambda \in \Lambda.
$$
Notice that the support of $p$ is contained in
$$
[0,1] + [0,1] + \dots + [0,1] = [0,n],
$$
and moreover, $p \in L^2[0,n]$ by Young's convolution theorem. Since $D(\Lambda) \geq n$, it follows from Fourier uniqueness that $p=0$. Applying the Fourier transform shows that
$
\ft p = \prod_{j=1}^n \ft p_j
$
vanishes identically. Since $p_j$ has compact support, each factor $\ft p_j$ is analytic. Hence, one of the factors must vanish identically which means that there exists $k \in \{0, \dots, n-1 \}$ so that $p_k = f- \omega_k h =0$. This shows that $E(\Lambda)$ does $\Theta$-PR in $L^2[0,1]$.

\textbf{Necessity.}
    Let $\Lambda$ be a lattice of density $D(\Lambda) < n$, i.e., $\Lambda = \alpha \Z$ for some $\alpha > n$. We prove that $E(\Lambda)$ does not do $\Theta$-PR in $L^2[0,1]$. In order to prove this, we show that the assumptions of Corollary \ref{cor:PW_recurrence} are satisfied.
    
    To do so, we define $\Lambda_j = \alpha(n\Z+j)$ for $j=0,...,n-1$, which implies that $\{ \Lambda_j \}_{j=0}^{n-1}$ is a cover of $\Lambda$. Next, we construct the functions $x_j$. To do so, let
    $$
    \xi \coloneqq \frac{1}{2n\alpha}.
    $$
    Since $2\leq n<\alpha < \infty$, it holds that $0 < \xi < \frac{1}{2}$. Further, define functions $S_j$ via
    $$
    S_j(x) = \sin \left ( 2\pi \xi x - \frac{j\pi}{n} \right ), \quad j \in \{ 0, \dots, n-1 \}.
    $$
    Clearly, each $S_j$ vanishes on $\Lambda_j$. Using the trigonometric identity
    $$
    \sin(v-v') + \sin(v+v') = 2 \sin(v)\cos(v'), \quad v,v' \in \R,
    $$
    and substituting $v = 2\pi \xi x  - (j-1)\frac{\pi}{n}$ and $v' = \frac{\pi}{n}$, it follows that
    $$
    S_j = 2 \cos(\tfrac{\pi}{n}) S_{j-1} - S_{j-2}, \quad j \geq 3
    $$
    Next, choose $\Phi \in PW$ such that $\hat \Phi$ has support in $[-(\frac{1}{2}-\xi),\frac{1}{2}-\xi]$, and $\Phi(0) \neq 0$. To obtain such a function, we can take $\phi \in L^2(\R) \setminus \{0 \}$ satisfying
    $$
    \supp(\phi) \subseteq [-(\tfrac{1}{2}-\xi),\tfrac{1}{2}-\xi], \quad \int_\R \phi(x) \, dx \neq 0,
    $$
    and define $\Phi \coloneqq \ift \phi$.
    Since the sine-function $S_j$ can be written as a linear combination of the exponentials $e^{2\pi i \xi x}$ and $e^{-2\pi i \xi x}$,
    $$
    S_j(x) = C e^{2\pi i \xi x} + C'e^{-2\pi i \xi x}, \quad C=-\frac{i}{2}\, e^{ -\frac{i j \pi}{n}}, \quad C' = \frac{i}{2}\, e^{\frac{i j \pi}{n}},
    $$
    it follows that multiplying $\Phi$ with $S_j$ shifts the support of $\mathcal{F}(S_j \Phi)$ by $\pm \xi$. This and the choice of $\Phi$ implies that $S_j \Phi$ is an element of $PW$. It also holds that $S_j \Phi$ is not the zero function: if it would be zero, then the product of the two analytic functions $C e^{2\pi i \xi x} + C'e^{-2\pi i \xi x}$ and $\Phi$ would vanish identically. Since $\Phi \neq 0$, it follows that $C e^{2\pi i \xi x} + C'e^{-2\pi i \xi x}$ is the zero function. Using the linear independence of complex exponentials, we have $C=C'=0$, which gives a contradiction.
    
Now let $\omega \coloneqq e^{\frac{2\pi i }{n}}$, $\zeta \coloneqq e^{\frac{2\pi i }{2n}}$, and define $x_j \in PW$ via
    $$
    x_j \coloneqq \zeta^j S_j \Phi.
    $$
    Then $x_j \neq 0$ and $x_j$ vanishes on $\Lambda_j$.
    Using the recurrence relation for $S_j$ it follows from a direct calculation that
    \begin{equation}\label{eq:Tj_relation}
        x_j = 2 \cos(\tfrac{\pi}{n}) \zeta \cdot  x_{j-1} - \zeta^2 \cdot  x_{j-2}.
    \end{equation}
    By Euler's identity we have
    $$
    2 \cos(\tfrac{\pi}{n}) \zeta = 2 \frac{e^{i\frac{\pi}{n}} + e^{-i\frac{\pi}{n}}}{2} e^{\frac{2\pi i }{2n}} = 1+\omega.
    $$
    Combining this with $\zeta^2 = \omega$, shows that \eqref{eq:Tj_relation} can be re-written as
    $$
    x_j = (1+\omega) x_{j-1} - \omega x_{j-2}.
    $$
\end{proof}

\end{document}
